\documentclass[12pt]{amsart}


\newtheorem{theorem}{Theorem}
\newtheorem{corollary}[theorem]{Corollary}
\newtheorem{lemma}[theorem]{Lemma}
\newcommand{\Z}{\mathbb{Z}}

\begin{document}
For integers $n\geq 1$, we define Gaussian integers $\Z[i]$
\[
P_n:=\prod_{k=1}^{n}(1+ik),
\]
where $A_n=\Re(P_n)$ and $B_n=\Im(P_n)$ are the real and imaginary parts respectively. Then define
$$
x_n:=\frac{B_n}{A_n}.
$$
The sequence begins with
\[
x_1=1,\ x_2=-3,\ x_3=0,\ x_4=4,\ x_5= -\tfrac{9}{19},\ x_6= \tfrac{105}{73},\ \dots.
\]
Define $\omega_n$ by
$$
\omega_n:=A_n^2+B_n^2=|P_n|^2.
$$
We define the \emph{squarefree kernel}
$$
K_n\;=\!\!\prod_{\nu_p(\omega_n)\ \mathrm{odd}}\!\! p,
$$
the product of primes dividing $\omega_n$ to an odd power.

\begin{theorem}\label{thm:main}
If $x_n=m$ is an integer, then the following hold.

\smallskip
\noindent
{\rm(1)} We have that $K_n\mid(1+m^2)$.

\smallskip
\noindent
{\rm(2)} If $K_n>1$, then
\[
|m|\ge\sqrt{K_n-1}.
\]
\end{theorem}

To isolate the indices at which $x_n$ might be an integer, we let
\[
E\;=\;\bigl\{\,n\ge5:\ |x_n|>\tfrac n2+1\,\bigr\}.
\]
By definition, if $n\notin E$, then $|x_n|\le n/2+1$.  For all
sufficiently large such $n$, this contradicts the lower bound
$|m|\ge\sqrt{K_n-1}$, and hence $x_n$ cannot be an integer. 
 We must count the exceptional set $E$. 
 
 \begin{lemma}\label{lem:proximity} If we let
\[
a_n:=\sum_{k=1}^{n}\arctan(1/k),
\]
then we have
\[
n\in E \quad\Longrightarrow\quad \exists j\in\Z\ \text{ such that }\
\left|a_n-j\frac{\pi}{2}\right|<\frac{2}{n}.
\]
\end{lemma}

\begin{corollary}\label{cor:density}
We have that
\[
\#(E\cap[1,N])=O(\log N).
\]
In particular, the conjecture holds on a set of density one.
\end{corollary}


\end{document}

